\textbf{The liquid is a reduced-order model that we wrote.} The ground truth is two modes with hand-chosen nonlinear terms, not a Navier--Stokes or particle simulation and not a real liquid. A learned model facing real sloshing meets breaking waves, three-dimensional modes and splashing; a pendulum model meets the same. Whether the ordering reported here survives is unknown. The nonlinearity strengths decide how wrong a single-mode model is, and we tested one setting plus the single-linear-mode extreme; we did not test a two-linear-mode truth, which would separate the effect of the second mode from that of the nonlinearity.

\textbf{Scale.} One planar axis, one container geometry, one transport distance, 5 seeds for the main table and 3 for ablations and sweeps, 50 episodes per run. Spill rates near zero are resolved only coarsely. The brief asked for 100 episodes per condition and seed and for a recurrent encoder; we used 250 episodes per condition pooled over seeds and a feedforward encoder to stay within the compute budget.

\textbf{The optimiser is not time-optimal.} The oracle MPC is slower than tuned open-loop input shaping on the nominal liquid, and the controller ablation moves the transport time by as much as the model choice does. Differences of a few hundredths of a second between MPC systems should not be read as properties of the models alone. A gradient-based or exact time-optimal solver would give a cleaner comparison.

\textbf{Tuning rule and grid.} Each system has one tuned scalar chosen on a grid from 90 validation episodes of the nominal liquid. The pendulum and spring-mass models, which predict almost equally well, received different margins from this rule; part of the reported gap between the learned and the pendulum controller is therefore a selection effect. Margins were not re-tuned for held-out liquids.

\textbf{Baselines.} All baselines are our reimplementations of textbook forms. The parametric models have fixed parameters identified on the nominal liquid; an adaptive pendulum with online parameter estimation, a robust input shaper, or a learned residual on a pendulum model would be stronger baselines on held-out liquids and were not tested. The learned model saw a range of liquids in training while the parametric models saw one, which favours the learned model on held-out liquids and the parametric models on the nominal liquid.

\textbf{Sensing.} The sensors are simulated proxies with noise levels we assigned; the camera is a noisy, delayed tilt signal, not an image pipeline. The learned model is trained with true surface heights as targets, which a real system would have to obtain from instrumentation. Grasp selection, arm kinematics, tilting the container and spill detection with stopping, all part of the original seed, are outside this study.
